\chapter{1977年甘肃卷}

\section{解答题}

\begin{problem}
用 \(>\)，\(<\)，\(=\) 号分别把下列各组数连接起来：

\begin{enumerate}
    \item \( 0.25+3\frac{1}{8} \) 与 \( \left| -3.375 \right| \)；
    \item \( \left( \frac{\sqrt{3}}{3}-\frac{1}{\sqrt{3}} \right)^3 \) 与 \( \left( -3 \right)^0 \)；
    \item \( -\lg 0.01 \) 与 \( \log_4 \frac{1}{16} \)；
    \item \( \tan 44^\circ \) 与 \( \cot 44^\circ \).
\end{enumerate}
\end{problem}

\begin{solution}
\begin{enumerate}
    \item \(0.25=\frac14\)，且 \(3\frac18=\frac{25}{8}\)，所以
    \(0.25+3\frac18=\frac14+\frac{25}{8}=\frac{27}{8}=3.375\)；又
    \(\left|-3.375\right|=3.375\). 因而
    \(0.25+3\frac18=\left|-3.375\right|\).

    \item 因为 \(\frac{1}{\sqrt3}=\frac{\sqrt3}{3}\)，所以
    \(\left(\frac{\sqrt3}{3}-\frac{1}{\sqrt3}\right)^3=0^3=0\)；又
    \(\left(-3\right)^0=1\). 因此
    \(\left(\frac{\sqrt3}{3}-\frac{1}{\sqrt3}\right)^3<\left(-3\right)^0\).

    \item \(0.01=10^{-2}\)，故 \(-\lg0.01=-(-2)=2\). 又
    \(\frac{1}{16}=4^{-2}\)，故 \(\log_4\frac{1}{16}=-2\). 因此
    \(-\lg0.01>\log_4\frac{1}{16}\).

    \item 在第一象限内，\(\cot44^\circ=\tan\left(90^\circ-44^\circ\right)=\tan46^\circ\). 正切函数在第一象限内单调递增，且 \(44^\circ<46^\circ\)，所以
    \(\tan44^\circ<\tan46^\circ=\cot44^\circ\).
\end{enumerate}
\end{solution}


\begin{problem}
解答下列各题：

\begin{enumerate}
    \item 计算：\( \frac{2a}{a-3b}-\frac{6ab-24b^2}{a^2-7ab+12b^2} \).
    \item 化简：\( \frac{\left( a^{\frac{8}{5}}b^{-\frac{6}{5}} \right)^{-\frac{1}{2}}\sqrt[5]{a^4}}{\sqrt[5]{b^3}} \).
    \item 证明：\( \frac{2\cot x}{1+\cot^2 x}=\sin 2x \).
    \item 求函数 \( y=\frac{1}{2}x^2+3x+9 \) 的极值.
\end{enumerate}
\end{problem}

\begin{solution}
\begin{enumerate}
    \item 分母不为零要求
    \(a-3b\ne0\) 且
    \(a^2-7ab+12b^2=(a-3b)(a-4b)\ne0\)，即 \(a\ne3b\) 且 \(a\ne4b\). 在此条件下，
    \[
    \text{原式}=\frac{2a}{a-3b}-\frac{6b(a-4b)}{(a-3b)(a-4b)}
    =\frac{2a}{a-3b}-\frac{6b}{a-3b}
    =\frac{2(a-3b)}{a-3b}=2
    \]

    \item 按分数指数幂的通常定义，题式有意义时 \(a>0\)，\(b>0\). 由
    \(\sqrt[5]{a^4}=a^{\frac45}\)、\(\sqrt[5]{b^3}=b^{\frac35}\)，并使用正数幂的运算法则，得
    \[
    \frac{\left( a^{\frac85}b^{-\frac65} \right)^{-\frac12}a^{\frac45}}{b^{\frac35}}
    =\frac{a^{-\frac45}b^{\frac35}a^{\frac45}}{b^{\frac35}}=1
    \]

    \item 在等式有意义的范围内，\(\sin x\ne0\)，并且
    \(1+\cot^2x=1+\frac{\cos^2x}{\sin^2x}=\frac{\sin^2x+\cos^2x}{\sin^2x}=\frac{1}{\sin^2x}\). 因而
    \[
    \frac{2\cot x}{1+\cot^2x}
    =\frac{2\frac{\cos x}{\sin x}}{\frac{1}{\sin^2x}}
    =2\sin x\cos x=\sin2x
    \]
    所以 \(\frac{2\cot x}{1+\cot^2x}=\sin2x\).

    \item 配方得
    \[
    y=\frac12x^2+3x+9=\frac12\left(x^2+6x\right)+9=\frac12\left(x+3\right)^2+\frac92
    \]
    因为 \(\left(x+3\right)^2\geqslant0\)，所以当 \(x=-3\) 时，\(y\) 取得最小值 \(\frac92\). 该二次函数开口向上，在 \(\R\) 上没有最大值.
\end{enumerate}
\end{solution}


\begin{problem}
\(AD\) 是 \( \triangle ABC \) 外接圆的直径，\( AD=6 \)，\( \angle DAC=\angle ABC \)，求 \(AC\)（精确到 \(0.01\)）.

\bitmapfigure[width=0.36\linewidth]{img/1977/gansu/q3_fig1.png}
\end{problem}

\begin{solution}
连接 \(CD\). 因为 \(AD\) 是圆的直径，所以由直径所对的圆周角为直角，得 \(\angle ACD=90^\circ\). 又因为 \(\angle ADC\) 与 \(\angle ABC\) 同对弦 \(AC\)，所以 \(\angle ADC=\angle ABC\). 结合已知 \(\angle ABC=\angle DAC\)，得到 \(\angle ADC=\angle DAC\). 因此 \(\triangle ADC\) 是等腰直角三角形，斜边 \(AD=6\)，故
\[
AC=AD\sin45^\circ=6\cdot\frac{\sqrt2}{2}=3\sqrt2\approx4.2426
\]
所以 \(AC\approx4.24\).
\end{solution}


\begin{problem}
在 \( \triangle ABC \) 中，\( \angle A=30^\circ \)，\( \angle B-\angle C=60^\circ \)，\( BC=2 \)，求 \(AC\)（用根式表示）.

\bitmapfigure[max width=0.52\linewidth]{img/1977/gansu/q4_fig1.png}
\end{problem}

\begin{solution}
三角形内角和给出 \(\angle B+\angle C=180^\circ-30^\circ=150^\circ\). 联立
\(\angle B+\angle C=150^\circ\) 与 \(\angle B-\angle C=60^\circ\)，得
\(\angle B=105^\circ\)，\(\angle C=45^\circ\). 由正弦定理，
\(\frac{AC}{\sin B}=\frac{BC}{\sin A}\)，\(AC=\frac{2\sin105^\circ}{\sin30^\circ}=4\sin105^\circ\).
再由和角公式，
\[
\sin105^\circ=\sin\left(60^\circ+45^\circ\right)
 =\sin60^\circ\cos45^\circ+\cos60^\circ\sin45^\circ
 =\frac{\sqrt6+\sqrt2}{4}
\]
所以 \(AC=\sqrt6+\sqrt2\).
\end{solution}


\begin{problem}
某生产队现已平整土地 \(100\text{ 亩}\)，计划两年后平整土地的总面积要在 \(225\text{ 亩}\) 以上，求每年平均增长率.
\end{problem}

\begin{solution}
原题没有说明“现已平整土地 \(100\text{ 亩}\)”是本年内新平整的面积还是截至目前的累计面积，也没有明确“两年后平整土地的总面积”是两年后的年平整面积还是累计面积，因此严格按原题不能唯一确定增长率.

若按通常的意图，将今年的平整面积作为基数，并将“两年后”理解为两年后的年平整面积，设每年平均增长率为 \(x\)，则 \(x\geqslant0\). 按原题解答给出的修订读法，并将“达到 \(225\text{ 亩}\) 以上”按严格超过理解，则
\[
100\left(1+x\right)^2>225
\]
于是 \(1+x>\frac32\)，从而 \(x>\frac12=50\%\). 若按通常“不小于”的语义理解“以上”，则不等号应为 \(\geqslant\)，得到 \(x\geqslant50\%\). 因此原题的文字不能唯一确定增长率；按原题解答给出的修订读法，结论为每年平均增长率大于 \(50\%\).
\end{solution}


\begin{problem}
右图的二视图表示什么形体？用直径为 \(40\text{ 毫米}\) 的圆柱形钢材锻压这样一个工件，需截取多长一段钢材？

\bitmapfigure[max width=0.44\linewidth]{img/1977/gansu/q6_fig1.png}
\end{problem}

\begin{solution}
由二视图可知，这个工件是一个圆台. 轴向截面中，上、下底面直径分别为 \(20\text{ 毫米}\) 和 \(60\text{ 毫米}\)，所以两底面半径分别为
\(r_1=10\text{ 毫米}\)、\(r_2=30\text{ 毫米}\). 由图中母线与底面成 \(45^\circ\)，且两底面半径之差为 \(20\text{ 毫米}\)，圆台高为
\[
h=\left(r_2-r_1\right)\tan45^\circ=20\text{ 毫米}
\]
圆台体积为
\[
V=\frac13\pi h\left(r_1^2+r_1r_2+r_2^2\right)
 =\frac13\pi\times20\left(10^2+10\times30+30^2\right)
 =\frac{26000\pi}{3}\text{ 毫米}^3
\]
设需截取长度为 \(l\text{ 毫米}\) 的圆柱形钢材. 锻压过程中体积不变，故
\[
\pi\left(\frac{40}{2}\right)^2l=\frac{26000\pi}{3}
\]
解得
\[
l=\frac{26000}{3\times20^2}=\frac{65}{3}\text{ 毫米}\approx21.7\text{ 毫米}
\]
所以需截取约 \(21.7\text{ 毫米}\) 长的钢材.
\end{solution}


\begin{problem}
一条直线通过 \( \left( 3,0 \right) \)，并且和直线 \( y=-2x+9 \) 垂直，求这条直线与直线 \( 2y+5x=3 \) 的交点的坐标.
\end{problem}

\begin{solution}
直线 \(y=-2x+9\) 的斜率为 \(-2\)，所以与它垂直的直线斜率为 \(\frac12\). 所求直线经过 \((3,0)\)，故其方程为
\(y-0=\frac12\left(x-3\right)\)，\(x-2y=3\).
设交点为 \((x,y)\)，则它满足
\[
\begin{cases}
 x-2y=3,\\
 2y+5x=3.
\end{cases}
\]
两式相加得 \(6x=6\)，所以 \(x=1\). 代入 \(x-2y=3\)，得 \(1-2y=3\)，即 \(y=-1\). 因此交点坐标为 \((1,-1)\).
\end{solution}


\begin{problem}
\begin{enumerate}
    \item 已知 \( \frac{dy}{dx}=x\e^x+2\cos x+3x^2 \)，求 \( y \).
    \item 过点 \( P\left( 4,0 \right) \) 向曲线 \( x^2+y^2=1 \) 作切线，求切线方程.
\end{enumerate}
\end{problem}

\begin{solution}
\begin{enumerate}
    \item 对等式两边关于 \(x\) 积分，得
    \[
    y=\int\left(x\e^x+2\cos x+3x^2\right)\,\mathrm dx
    \]
    分别计算各项：分部积分给出
    \(\int x\e^x\,\mathrm dx=x\e^x-\int\e^x\,\mathrm dx=\e^x(x-1)\)，另有
    \(\int2\cos x\,\mathrm dx=2\sin x\)、\(\int3x^2\,\mathrm dx=x^3\). 因此
    \[
    y=\e^x(x-1)+2\sin x+x^3+C
    \]
    其中 \(C\) 是任意常数.

    \item 设切点为 \(T\left(x_0,y_0\right)\). 因为圆心为原点，半径 \(OT\) 垂直于切线，所以切线方程为
    \[
    x_0x+y_0y=x_0^2+y_0^2=1
    \]
    又因为 \(T\) 在圆 \(x^2+y^2=1\) 上，所以 \(x_0^2+y_0^2=1\). 点 \(P(4,0)\) 在切线上，代入切线方程得 \(4x_0=1\)，即 \(x_0=\frac14\). 于是
\(y_0^2=1-\left(\frac14\right)^2=\frac{15}{16}\)，\(y_0=\pm\frac{\sqrt{15}}4\).
    两个切点为 \(\left(\frac14,\frac{\sqrt{15}}4\right)\) 和 \(\left(\frac14,-\frac{\sqrt{15}}4\right)\). 分别代入 \(x_0x+y_0y=1\)，得到两条切线
\(x+\sqrt{15}y=4\)，\(x-\sqrt{15}y=4\).
\end{enumerate}
\end{solution}
